\uuid{Y4JR}
\titre{Intervalle}
\auteur{bodin ; hanani ; mzari}
\organisation{exo7}
\licence{CC BY-NC-SA 4.0}
\datecreate{2020-01-20}
\niveau{L1}
\module{Analyse}
\chapitre{Propriétés de R}
\sousChapitre{Intervalle, densité}
\difficulte{2}
\theme{Intervalle, densité, intervalle}
\exo7qcmid{lille-1.reels.024}

\contenu{
\begin{question}
\qclassification{120.04}
\qtags{motcle=intervalle}
Pour $n\ge 1$ on définit l'intervalle $I_n = [0,\frac1n]$.
Quelles sont les assertions vraies ?
\begin{answers}
    \bad{Pour tout $n\ge 1$, $I_n \subset I_{n+1}$.}

    \good{Si $x \in I_n$ pour tout $n\ge1$, alors $x=0$.}

    \bad{L'union de tous les $I_n$ (pour $n$ parcourant $\Nn^*$) est $[0,+\infty[$.}

    \bad{Pour $n < m$ alors $I_n \cap I_{n+1} \cap \ldots \cap I_m = I_n$.}
\end{answers}
\begin{explanations}
On a $[0,1] = I_1 \supset I_2 \supset I_3 \supset \cdots$.
\end{explanations}
\end{question}
}
